\section{Software, Tools and Applications}

\paragraph{Languages and build tools.}
Google's build tool Blaze, open sourced as Bazel, builds binaries from what Google calls its standard languages: C++, Java, Python, Go and JavaScript \cite{sre-book}. Builds are hermetic, meaning they depend on known versions of tools such as compilers and not on software installed on the build machine \cite{sre-book}. On scripting, the SRE book describes the shell scripts Google first used to configure clusters as brittle and says the initial automation was simple Python scripts \cite{sre-book}.

\paragraph{Editors, compilers and IDEs.}
Compiler plug-ins such as Error Prone for Java and Tsetse for TypeScript can prohibit risky code patterns at compile time, and compiler errors give faster feedback than linters or checks at code review time \cite{bsrs}. We found no book that names an editor or an IDE as an SRE tool.

\paragraph{Debugging and profiling.}
Every Google server has an HTTP server that provides diagnostics and statistics for debugging \cite{sre-book}. Valgrind catches memory errors without recompiling, AddressSanitizer detects buffer overflows and use after free, and ThreadSanitizer detects data races \cite{bsrs}. These claims do not show that Google SREs use these tools every day.

\paragraph{Version control, containers and shell.}
The SRE book says that Google stores all of its schemes for distributing configuration in the primary source code repository and enforces a strict code review requirement \cite{sre-book}. The SRE Workbook calls checking configuration files into a versioning system such as Subversion or Git common practice \cite{workbook}. It describes Kubernetes as an open source system for orchestrating containerized workloads on a cluster of machines, used by writing YAML files and running kubectl \cite{workbook}. The SRE book calls Kubernetes Borg's descendant, started by Google in 2014 \cite{sre-book}. Much of Google's distributed automation once relied on SSH, which was clumsy from a security perspective because people needed root on many machines, so Google replaced sshd with an authenticated, ACL-driven, RPC-based Local Admin Daemon \cite{sre-book}. \emph{Building Secure and Reliable Systems} says that, instead of setting services up by hand, you might write imperative turnup automation or use a declarative configuration language such as Terraform \cite{bsrs}.

\paragraph{Release engineering and CI/CD.}
The SRE book says that Google has an automated release system called Rapid, that some teams use a ``Push on Green'' model in which they deploy every build that passes all tests, and that a continuous test system runs unit tests against the mainline each time a change is submitted \cite{sre-book}. The Workbook defines canarying as a partial and time-limited deployment of a change and its evaluation \cite{workbook}. \emph{Building Secure and Reliable Systems} describes the software supply chain as writing, building, testing and deploying a system through version control, a continuous integration (CI) pipeline and a continuous delivery (CD) pipeline, and names Jenkins as an example \cite{bsrs}. In one Workbook case study, a small group at Waze, which Google acquired in 2013, proposed continuous deployment with Spinnaker. This is one team's case and not general Google practice \cite{workbook}.

\paragraph{Monitoring, alerting and incident response.}
The SRE book says that Google built its monitoring system Borgmon shortly after it created the Borg job scheduler in 2003. Borgmon fetches the \texttt{/varz} URI of each monitored target at set intervals, and the idea of using time-series data to generate alerts is available outside Google through open source tools such as Prometheus \cite{sre-book}. The Workbook strongly recommends treating monitoring configuration as code \cite{workbook}. The SRE book names latency, traffic, errors and saturation as the four golden signals and Dapper as a tool for tracing requests \cite{sre-book}. Alertmanager routes alert notifications, page-worthy alerts go to the on-call rotation and subcritical ones go to ticket queues \cite{sre-book}. The Workbook names three main roles in incident response: the Incident Commander, the Communications Lead and the Operations Lead \cite{workbook}.

\paragraph{Server and database software.}
Google's Ads database, a MySQL deployment, ran on Borg on shared machines that were rebooted for kernel upgrades \cite{sre-book}. The report \textit{Creating a Production Launch Plan} gives the example of a relational database with an upper limit on transaction throughput, where adding CPU or memory scales only so far \cite{static-pdf-creatingproductionlaunchplan-pdf}. Caches sometimes become hard dependencies of a service \cite{sre-book}. A message queue that stores data only in volatile memory can lose transactions \cite{bsrs}, and in the Waze case study the message queueing system was critical because every component used it \cite{workbook}. This section does not cover web servers such as Apache or Nginx, or the operation of PostgreSQL, MongoDB or Redis.

